\documentclass[11pt]{article}

\usepackage[utf8]{inputenc}
\usepackage{amsmath, amssymb, amsthm}
\usepackage{geometry}
\usepackage{hyperref}
\usepackage{listings}
\usepackage{xcolor}

% Forzar que las palabras automáticas (como "Abstract") salgan en inglés
\usepackage[english]{babel}

\geometry{margin=1in}

% Configuración para bloques de código
\lstset{
    language=Python,
    basicstyle=\ttfamily\small,
    keywordstyle=\color{blue},
    stringstyle=\color{red},
    commentstyle=\color{gray},
    numbers=left,
    numberstyle=\tiny\color{gray},
    stepnumber=1,
    numbersep=8pt,
    frame=single,
    breaklines=true
}

% Definición correcta de operadores matemáticos
\DeclareMathOperator{\Cov}{Cov}

\title{Filters Can Also Create Correlations: A Reproducible Computational Model and Its Implications}

\author{German Garcia}
\date{2026}

\newtheorem{proposition}{Proposition}

\begin{document}

\maketitle

\begin{abstract}
We study how selection mechanisms can induce apparent correlations between originally independent variables. Using a minimal computational model, we show that conditioning on a selection region can generate non-zero covariance. We present a corrected formal proposition, emphasizing that induced correlation depends on the geometry of the selection region.
\end{abstract}

\section{Introduction}

Observed correlation in data may depend not only on underlying relationships but also on the selection process. This work shows how a filter can induce apparent dependence between independent variables. 

This phenomenon was originally described in the medical field by Berkson \cite{berkson1946} and later integrated into modern graph theory and causal analysis by Pearl \cite{pearl2000} under the term \textit{collider bias}.

\section{Model}

Let:
\[
X \sim \mathcal{N}(0,1), \quad Y \sim \mathcal{N}(0,1)
\]
be independent variables.

We define:
\[
S(X,Y) = X + Y
\]
and a selection region:
\[
\Omega = \{(x,y) \in \mathbb{R}^2 : x+y > 1.0\}
\]

\section{Computational Results}

Empirically:
\begin{itemize}
\item Before the filter: $\Cov(X,Y) \approx 0$
\item After the filter: $\Cov(X,Y \mid \Omega) \neq 0$ (in general)
\end{itemize}

\section{Corrected Proposition}

\begin{proposition}[Selection-induced dependence]
Let $X$ and $Y$ be independent random variables with a continuous joint density.

Then, there exist selection regions $\Omega \subset \mathbb{R}^2$ such that:
\[
\Cov(X,Y \mid \Omega) \neq 0
\]
even though:
\[
\Cov(X,Y) = 0
\]
\end{proposition}

\textbf{Proof.}

The original joint density factorizes:
\[
f_{X,Y}(x,y) = f_X(x)f_Y(y)
\]

By conditioning on $\Omega$, the density becomes:
\[
f_{X,Y|\Omega}(x,y) \propto f_X(x)f_Y(y)\mathbf{1}_\Omega(x,y)
\]

If $\Omega$ is not a Cartesian product (i.e., it cannot be written as $A \times B$), then the integration region introduces dependence.

For example, if:
\[
\Omega = \{(x,y): x+y>1.0\}
\]
then high values of $x$ restrict the possible values of $y$, generating a correlation.

Therefore, there exist regions $\Omega$ such that:
\[
\Cov(X,Y|\Omega) \neq 0
\]
\hfill $\square$

\section{Computational Simulation}

To empirically validate the proposition, the following Python code generates independent variables and demonstrates the induction of correlation after applying the selection filter:

\begin{lstlisting}
import numpy as np

# Seed for reproducibility
np.random.seed(42)

# Generation of independent variables
n_samples = 100000
X = np.random.normal(0, 1, n_samples)
Y = np.random.normal(0, 1, n_samples)

print(f"Original Covariance: {np.cov(X, Y):.4f}")

# Selection filter (e.g., truncate by the sum)
mask = (X + Y) > 1.0
X_filtered = X[mask]
Y_filtered = Y[mask]

print(f"Covariance after filter: {np.cov(X_filtered, Y_filtered):.4f}")
\end{lstlisting}

\section{Conclusion}

A selection filter can induce apparent correlations between independent variables, but the effect critically depends on the geometry of the selection region.

\begin{thebibliography}{9}

\bibitem{berkson1946}
Berkson, J. (1946). 
Limitations of the application of the contingency table analysis to hospital data. 
\textit{Biometrics Bulletin}, 2(3), 47-53.

\bibitem{pearl2000}
Pearl, J. (2000). 
\textit{Causality: Models, Reasoning and Inference}. 
Cambridge University Press.

\end{thebibliography}

\end{document}

